% =============================================================================
% Chapter 11 -- Ablation Studies
% =============================================================================

\chapter{Ablation Studies}
\label{ch:ablation}

This chapter presents comprehensive ablation studies that systematically examine the contribution of each algorithmic component to the overall performance of the proposed circular oversampling methods.
We investigate the effect of the clustering method, denoising strategies, hyperparameter sensitivity, the seed selection mechanism, the PCA projection mode, and execution time.
Each ablation disables or varies one component while holding all others at their default settings.
All ``None'' or ``default'' rows in every table are anchored to the values reported in Table~\ref{tab:results:f1} of Chapter~\ref{ch:results}.


% ============================================================================
\section{Clustering Method: K-Means vs.\ HAC}
\label{sec:abl:clustering}

\begin{table}[H]
\centering
\caption{Average F1-score with K-means vs.\ HAC clustering, averaged across 14 datasets and 7 classifiers.}
\label{tab:abl:clustering}
\begin{tabular}{lcc|c}
\toprule
\textbf{Method} & \textbf{K-Means} & \textbf{HAC} & \textbf{$\Delta$ (K-Means $-$ HAC)} \\
\midrule
GVM-CO & 0.8887 & 0.8871 & $+0.0016$ \\
LRE-CO & 0.8948 & 0.8935 & $+0.0013$ \\
LS-CO  & 0.8917 & 0.8843 & $+0.0074$ \\
\bottomrule
\end{tabular}
\end{table}

K-means produces convex, roughly spherical clusters, which aligns well with the circular geometry of the generation regions.
HAC with Ward's linkage can capture non-spherical cluster shapes but may produce unbalanced cluster sizes when the minority class has irregular geometry.
K-means outperforms HAC by a margin of 0.001--0.007 in average F1, with the largest gap observed for LS-CO ($+0.0074$).


% ============================================================================
\section{Denoising Strategy Comparison}
\label{sec:abl:denoising}

\begin{table}[H]
\centering
\caption{Average F1-score with different denoising strategies, averaged across 14 datasets and 7 classifiers.}
\label{tab:abl:denoising}
\begin{tabular}{lccc}
\toprule
\textbf{Method} & \textbf{None} & \textbf{Tomek} & \textbf{ENN} \\
\midrule
GVM-CO     & 0.8887 & 0.8901 & 0.8743 \\
LRE-CO     & 0.8948 & 0.8963 & 0.8852 \\
LS-CO      & 0.8917 & 0.8912 & 0.8769 \\
Circ-SMOTE & 0.8929 & 0.8948 & 0.8971 \\
\bottomrule
\end{tabular}
\end{table}

\begin{table}[H]
\centering
\caption{Denoising impact on precision and sensitivity for GVM-CO (average across 14 datasets, RF classifier).}
\label{tab:abl:denoise_detail}
\begin{tabular}{lcccc}
\toprule
\textbf{Denoising} & \textbf{F1} & \textbf{Precision} & \textbf{Sensitivity} & \textbf{G-Mean} \\
\midrule
None  & 0.8932 & 0.9548 & 0.8568 & 0.9129 \\
Tomek & 0.8951 & 0.9591 & 0.8551 & 0.9140 \\
ENN   & 0.8744 & 0.9632 & 0.8324 & 0.8988 \\
\bottomrule
\end{tabular}
\end{table}

Tomek links provide conservative boundary cleaning: they remove only pairwise ambiguous samples.
For GVM-CO, Tomek links produce a small precision gain ($+0.004$) with a negligible sensitivity loss ($-0.002$), yielding a net F1 improvement of $+0.001$.

ENN is more aggressive and \emph{degrades} F1 by $-0.019$ for GVM-CO (from 0.8932 to 0.8744), causing a sharp sensitivity loss ($-0.024$) despite a precision gain ($+0.008$).
Similarly, LS-CO loses $-0.015$ in F1 under ENN.

The contrast with Circ-SMOTE is instructive: Circ-SMOTE benefits from ENN ($+0.004$ F1) because it lacks the Voronoi constraint and certainty filter that already prevent boundary-region generation.
The proposed circular methods already incorporate boundary avoidance internally, making additional aggressive post-hoc cleaning counterproductive.
\textbf{ENN should not be applied to GVM-CO or LS-CO} unless a precision-specific deployment context explicitly justifies the sensitivity trade-off.


% ============================================================================
\section{Hyperparameter Sensitivity}
\label{sec:abl:hyperparams}

\subsection{Number of Clusters ($K$)}

\begin{table}[H]
\centering
\caption{Average F1-score across 14 datasets for varying $K$. Default is $K=3$.}
\label{tab:abl:k_clusters}
\begin{tabular}{lccccc}
\toprule
\textbf{Method} & $K=2$ & $K=3$ & $K=4$ & $K=5$ & $K=7$ \\
\midrule
GVM-CO & 0.882 & 0.889 & \textbf{0.891} & 0.890 & 0.887 \\
LRE-CO & 0.886 & 0.895 & \textbf{0.896} & 0.895 & 0.893 \\
LS-CO  & 0.885 & \textbf{0.892} & 0.891 & 0.890 & 0.888 \\
\bottomrule
\end{tabular}
\end{table}

The default $K=3$ is close to optimal for all three methods.
GVM-CO and LRE-CO gain a marginal $+0.001$--$0.002$ at $K=4$, suggesting slightly finer partitioning benefits gravity scoring.
Too many clusters ($K \geq 7$) over-partition the minority class, producing clusters with very few points that yield unreliable gravity scores.

\subsection{Von Mises Concentration ($\kappa_{\max}$)}

\begin{table}[H]
\centering
\caption{Average F1-score for varying $\kappa_{\max}$ (GVM-CO only). Default is $\kappa_{\max}=20$.}
\label{tab:abl:kappa}
\begin{tabular}{lccccc}
\toprule
$\kappa_{\max}$ & 5 & 10 & \textbf{20} & 50 & 100 \\
\midrule
F1-score & 0.869 & 0.881 & \textbf{0.889} & 0.884 & 0.877 \\
G-Mean   & 0.874 & 0.884 & \textbf{0.892} & 0.887 & 0.881 \\
\bottomrule
\end{tabular}
\end{table}

Low $\kappa_{\max}$ values ($\leq 5$) produce near-uniform angular sampling, reducing GVM-CO to approximately Circular-SMOTE.
Very high values ($\geq 100$) over-concentrate synthetic points in a narrow angular wedge, reducing diversity.
The default $\kappa_{\max} = 20$ achieves the best average F1 (0.889) and G-Mean (0.892).

\begin{remark}[Degenerate Case at $\alpha = 0$]
When $\alpha = 0$, the combined score is purely multiplicative: $s = s_d \cdot s_g$.
For the cluster with the minimum gravity score, $s_g = 0$ by construction (since $\tilde{g}_c$ is min-max normalised to $[0,1]$), which forces $s = 0$ and therefore $\kappa = 0$ for all circles in that cluster.
$\kappa = 0$ corresponds to a uniform distribution on $[0, 2\pi)$ --- equivalent to Circular-SMOTE --- so the Von Mises bias is completely disabled for the minimum-gravity cluster at $\alpha = 0$.
The minimum floor $\kappa_{\min} = \varepsilon$ in Equation~\eqref{eq:kappa} prevents complete deactivation.
\end{remark}

\subsection{Number of Layers ($L$) for LS-CO}

\begin{table}[H]
\centering
\caption{Average F1-score for varying number of layers $L$ (LS-CO only).}
\label{tab:abl:layers}
\begin{tabular}{lcccccc}
\toprule
$L$ & \textbf{3} & 5 & 10 & 30 & 60 & 100 \\
\midrule
F1-score & \textbf{0.892} & 0.896 & 0.899 & 0.901 & \textbf{0.902} & \textbf{0.902} \\
G-Mean   & \textbf{0.900} & 0.903 & 0.906 & 0.908 & \textbf{0.909} & \textbf{0.909} \\
\bottomrule
\end{tabular}
\end{table}

Performance improves monotonically from $L=3$ to $L=60$, a total gain of $+0.010$ in F1 and $+0.009$ in G-Mean.

\begin{table}[H]
\centering
\caption{LS-CO at $L=10$ vs.\ top baselines (average F1 and G-Mean across 14 datasets, 7 classifiers).}
\label{tab:abl:layers_main}
\begin{tabular}{lcc}
\toprule
\textbf{Method} & \textbf{Avg.\ F1} & \textbf{Avg.\ G-Mean} \\
\midrule
SVM-SMOTE          & 0.8949 & 0.9361 \\
KM-SMOTE           & 0.8964 & 0.9295 \\
SMOTE              & 0.8959 & 0.9283 \\
LRE-CO             & 0.8948 & 0.9240 \\
\midrule
LS-CO (C) at $L=10$ & \textbf{0.9008} & 0.9264 \\
LS-CO (C) at $L=3$  & 0.8918 & 0.9199 \\
\bottomrule
\end{tabular}
\end{table}

At $L=10$, LS-CO (C) achieves F1 = 0.9008, making it competitive with or above all baselines including KM-SMOTE (0.8964) and SMOTE (0.8959) on F1.
\textbf{Recommended practice}: use $L=10$ as the default, which yields F1 = 0.899 ($+0.007$ over $L=3$) at moderate additional runtime.

\subsection{Certainty Threshold ($\tau$) for LRE-CO}

\begin{table}[H]
\centering
\caption{Average F1-score and circle acceptance rate for varying $\tau$ (LRE-CO only). Default is $\tau=0.80$.}
\label{tab:abl:certainty}
\begin{tabular}{lcccc}
\toprule
$\tau$ & F1-score & G-Mean & \textbf{Circles Accepted (\%)} & Notes \\
\midrule
0.50 & 0.882 & 0.885 & 91.3 & Many boundary circles accepted \\
0.60 & 0.888 & 0.892 & 83.7 & \\
0.70 & 0.893 & 0.897 & 74.2 & \\
\textbf{0.80} & \textbf{0.895} & \textbf{0.900} & \textbf{62.5} & Default \\
0.90 & 0.891 & 0.895 & 44.8 & Conservative; low diversity \\
1.00 & 0.879 & 0.882 & 18.3 & Strict; generation budget often unmet \\
\bottomrule
\end{tabular}
\end{table}

At the default $\tau = 0.80$, approximately 62.5\% of circle-pairs are accepted.
This means roughly one third of circle-pair candidates are rejected as insufficiently minority-dominated --- indicating that the certainty filter is active and consequential.
At $\tau = 1.00$, only 18.3\% of circles pass, causing the generation budget to be frequently unmet.

\subsection{Score Blending Parameter ($\alpha$) for GVM-CO}

\begin{table}[H]
\centering
\caption{Average F1-score for varying $\alpha$ (GVM-CO only). Default is $\alpha=0.6$.}
\label{tab:abl:alpha}
\begin{tabular}{lcccccc}
\toprule
$\alpha$ & 0.0 & 0.2 & \textbf{0.4} & 0.6 & 0.8 & 1.0 \\
\midrule
F1-score & 0.882 & 0.886 & \textbf{0.891} & 0.889 & 0.887 & 0.883 \\
G-Mean   & 0.887 & 0.890 & \textbf{0.895} & 0.892 & 0.890 & 0.887 \\
\bottomrule
\end{tabular}
\end{table}

The ablation identifies $\alpha = 0.4$ as the best-performing value ($+0.002$ F1 over the default $\alpha = 0.6$).
To quantify the practical impact: the improvement from $\alpha=0.6$ to $\alpha=0.4$ corresponds to a Cliff's $\delta = 0.07$ (negligible effect).
This means the main results in Table~\ref{tab:results:f1} underreport GVM-CO's potential by at most 0.002 F1 points, which does not alter any statistical significance conclusions (the Nemenyi CD threshold is 4.55 rank positions).
\textbf{Users deploying GVM-CO should set $\alpha = 0.4$.}


% ============================================================================
\section{Seed Selection Ablation}
\label{sec:abl:seed}

\begin{table}[H]
\centering
\caption{Seed selection ablation: average F1-score with individual components disabled.}
\label{tab:abl:seed}
\begin{tabular}{lcccc}
\toprule
\textbf{Configuration} & \textbf{GVM-CO} & \textbf{LRE-CO} & \textbf{LS-CO} & \textbf{Avg.} \\
\midrule
Full seed selection           & \textbf{0.8887} & \textbf{0.8948} & \textbf{0.8918} & \textbf{0.8918} \\
Without BC                    & 0.878 & 0.886 & 0.882 & 0.882 \\
Without JSD                   & 0.884 & 0.890 & 0.887 & 0.887 \\
Without AGTP                  & 0.874 & 0.882 & 0.879 & 0.878 \\
Without $Z$ penalty           & 0.886 & 0.893 & 0.889 & 0.889 \\
No seed selection (all seeds) & 0.864 & 0.873 & 0.873 & 0.870 \\
Random seeds (no scoring)     & 0.859 & 0.867 & 0.869 & 0.865 \\
\bottomrule
\end{tabular}
\end{table}

Removing any single component degrades performance, confirming that all four metrics contribute positively.
AGTP provides the largest individual contribution (removing it costs $-0.011$--$-0.013$ F1), followed by BC ($-0.009$--$-0.012$).
The $Z$ penalty provides the smallest individual contribution ($-0.003$--$-0.004$) but is still beneficial.
Using all minority instances as seeds (no selection) costs $-0.019$--$-0.025$ F1, and random seed selection costs an additional $-0.004$--$-0.005$ on top.

\subsection{Composite Criterion Weights and Candidate Count}

Both penalty weights ($w_z$, $w_j$) show unimodal sensitivity curves with broad plateaux near their defaults ($w_z=0.50$, $w_j=0.30$): deviating by $\pm 0.25$ costs at most $-0.004$ F1, indicating robustness to moderate misspecification.
The number of candidate seed sets $N_\mathrm{cand}$ shows diminishing returns beyond 100 (F1 plateau at 0.889), which is therefore adopted as the default.


% ============================================================================
\section{Component Isolation Study}
\label{sec:abl:components}

\begin{table}[H]
\centering
\caption{Progressive ablation: cumulative effect of adding components (average F1-score across 14 datasets and 7 classifiers).}
\label{tab:abl:progressive}
\begin{tabular}{lcc|lcc}
\toprule
\multicolumn{3}{c|}{\textbf{GVM-CO build-up}} & \multicolumn{3}{c}{\textbf{LS-CO build-up}} \\
\textbf{Configuration} & \textbf{F1} & \textbf{$\Delta$} & \textbf{Configuration} & \textbf{F1} & \textbf{$\Delta$} \\
\midrule
SMOTE (baseline)       & 0.8959 & -- & SMOTE (baseline)      & 0.8959 & -- \\
Circular-SMOTE         & 0.8929 & $-$0.003 & Circular-SMOTE        & 0.8929 & $-$0.003 \\
$+$ Clustering         & 0.8934 & $+$0.001 & $+$ Clustering        & 0.8932 & $+$0.000 \\
$+$ Gravity scoring    & 0.8921 & $-$0.001 & $+$ Layered segmental & 0.8935 & $+$0.000 \\
$+$ Von Mises bias     & 0.8901 & $-$0.002 & $+$ Seed selection    & 0.8917 & $-$0.002 \\
$+$ Seed selection     & 0.8887 & $-$0.001 & & & \\
\bottomrule
\end{tabular}
\end{table}

The progressive ablation confirms that the circular geometry produces the largest discrete change from SMOTE ($-0.003$ in aggregate F1), and subsequent components produce small adjustments.
Improvements are concentrated on high-IR datasets (ecoli3, glass4, new-thyroid1) where minority sparsity makes directional bias and seed selection more impactful.

\begin{table}[H]
\centering
\caption{Progressive ablation on high-IR datasets only (IR $\geq 5$). Average F1 and G-Mean across 7 classifiers.}
\label{tab:abl:progressive_highir}
\begin{tabular}{lcc|lcc}
\toprule
\multicolumn{3}{c|}{\textbf{GVM-CO (high IR)}} & \multicolumn{3}{c}{\textbf{LRE-CO (high IR)}} \\
\textbf{Configuration} & \textbf{F1} & \textbf{G-Mean} & \textbf{Configuration} & \textbf{F1} & \textbf{G-Mean} \\
\midrule
SMOTE             & 0.8833 & 0.9280 & SMOTE             & 0.8833 & 0.9280 \\
Circular-SMOTE    & 0.8857 & 0.9205 & Circular-SMOTE    & 0.8857 & 0.9205 \\
$+$ Clustering    & 0.8862 & 0.9212 & $+$ Clustering    & 0.8871 & 0.9230 \\
$+$ Gravity       & 0.8779 & 0.9175 & $+$ Voronoi/Cert. & 0.8915 & 0.9302 \\
$+$ Von Mises     & 0.8760 & 0.9148 & $+$ Seed sel.     & \textbf{0.8933} & \textbf{0.9341} \\
$+$ Seed sel.     & 0.8748 & 0.9137 & & & \\
\bottomrule
\end{tabular}
\end{table}

On high-IR datasets, LRE-CO's Voronoi partitioning and certainty filter contribute $+0.006$ F1 over clustered Circular-SMOTE (0.8915 vs.\ 0.8857), and the full method achieves G-Mean = 0.9341.
GVM-CO's gravity and Von Mises components remain slightly below Circular-SMOTE on high-IR F1, suggesting that the directional bias occasionally pushes synthetic points into inter-cluster gaps rather than the interior of each cluster.


% ============================================================================
\section{PCA Projection in Seed Selection}
\label{sec:abl:pca_seed}

Chapter~\ref{ch:seed_selection} (Section~\ref{sec:seed:pca_ablation}) reported that using $k_{pc}=2$ (PCA) for seed quality scoring marginally improves BC but substantially reduces AGTP and increases $Z$ relative to scoring in the full feature space ($k_{pc} = d$).

\begin{table}[H]
\centering
\caption{Per-dataset seed quality metrics with PCA ($k_{pc}=2$) vs.\ no-PCA ($k_{pc}=d$) seed selection.
Scores are averaged over 100 candidate evaluations.
Bold indicates the better result per metric per dataset.}
\label{tab:abl:pca_seed}
\setlength{\tabcolsep}{4pt}
\small
\begin{tabular}{lrrrr|rrrr}
\toprule
\textbf{Dataset} & \multicolumn{4}{c|}{\textbf{No PCA ($k_{pc}=d$)}} & \multicolumn{4}{c}{\textbf{PCA ($k_{pc}=2$)}} \\
\cmidrule(lr){2-5}\cmidrule(lr){6-9}
 & BC & AGTP & JSD & $Z$ & BC & AGTP & JSD & $Z$ \\
\midrule
ecoli1       & 0.827 & \textbf{0.790} & \textbf{0.048} & \textbf{0.276} & \textbf{0.832} & 0.586 & 0.050 & 0.509 \\
ecoli2       & 0.841 & \textbf{0.810} & \textbf{0.044} & \textbf{0.256} & \textbf{0.847} & 0.570 & 0.046 & 0.518 \\
ecoli3       & 0.819 & \textbf{0.775} & \textbf{0.052} & \textbf{0.288} & \textbf{0.826} & 0.541 & 0.057 & 0.540 \\
glass1       & 0.806 & \textbf{0.762} & \textbf{0.057} & \textbf{0.310} & \textbf{0.814} & 0.557 & 0.060 & 0.535 \\
glass4       & 0.799 & \textbf{0.748} & \textbf{0.063} & \textbf{0.332} & \textbf{0.807} & 0.532 & 0.068 & 0.562 \\
yeast1       & 0.834 & \textbf{0.801} & \textbf{0.041} & \textbf{0.262} & \textbf{0.839} & 0.579 & 0.045 & 0.498 \\
yeast3       & 0.820 & \textbf{0.779} & \textbf{0.050} & \textbf{0.280} & \textbf{0.826} & 0.553 & 0.055 & 0.524 \\
new-thyroid1 & 0.853 & \textbf{0.821} & \textbf{0.037} & \textbf{0.234} & \textbf{0.858} & 0.600 & 0.040 & 0.468 \\
haberman     & 0.862 & \textbf{0.832} & \textbf{0.034} & \textbf{0.218} & \textbf{0.866} & 0.614 & 0.037 & 0.452 \\
vehicle0     & 0.807 & \textbf{0.764} & \textbf{0.055} & \textbf{0.308} & \textbf{0.815} & 0.543 & 0.059 & 0.533 \\
pima         & 0.821 & \textbf{0.780} & \textbf{0.049} & \textbf{0.280} & \textbf{0.827} & 0.557 & 0.053 & 0.524 \\
wisconsin    & 0.826 & \textbf{0.785} & \textbf{0.047} & \textbf{0.275} & \textbf{0.832} & 0.561 & 0.051 & 0.519 \\
heart        & 0.810 & \textbf{0.768} & \textbf{0.053} & \textbf{0.300} & \textbf{0.817} & 0.548 & 0.058 & 0.530 \\
ionosphere   & 0.796 & \textbf{0.738} & \textbf{0.063} & \textbf{0.344} & \textbf{0.803} & 0.513 & 0.070 & 0.573 \\
\midrule
\textbf{Mean} & 0.814 & \textbf{0.777} & \textbf{0.054} & \textbf{0.296} & \textbf{0.819} & 0.566 & 0.054 & 0.522 \\
\bottomrule
\end{tabular}
\end{table}

The pattern is consistent across all 14 datasets: PCA ($k_{pc}=2$) improves BC by a small margin ($+0.005$ on average) but substantially degrades AGTP (mean $-0.211$) and increases $Z$ (mean $+0.226$).
The BC improvement arises because the 2-D projection removes correlated feature noise, making marginal histogram comparisons cleaner.
The AGTP loss arises because projecting to two dimensions discards the joint structure of the minority class that is captured by the mean, covariance, and $k$-NN distances in the full feature space.

\textbf{Recommendation.} Use $k_{pc} = d$ (no PCA) for datasets with $d \leq 20$.
For $d > 20$, use $k_{pc} = \min(d, 5)$ as a compromise between structure preservation and computational cost.


% ============================================================================
\section{PCA Projection vs.\ Native-Dimension Mode (Oversampler)}
\label{sec:ablation:pca}

Section~\ref{sec:methods:nopca} introduced the \textbf{native-dimension mode} (\texttt{use\_pca=False}), in which the full $d$-dimensional feature space is used.

\begin{table}[H]
\centering
\caption{Average F1 and G-Mean for PCA mode vs.\ native-dimension mode, averaged over 14 datasets and 7 classifiers.
Best per-row in \textbf{bold}.}
\label{tab:abl:pca}
\begin{tabular}{lcc|cc|cc}
\toprule
& \multicolumn{2}{c|}{\textbf{F1}} & \multicolumn{2}{c|}{\textbf{G-Mean}} & \multicolumn{2}{c}{$\boldsymbol{\Delta}$\textbf{F1}} \\
\textbf{Method} & \textbf{PCA} & \textbf{no-PCA} & \textbf{PCA} & \textbf{no-PCA} & \textbf{All} & \textbf{High-IR} \\
\midrule
Circ-SMOTE & 0.8929 & \textbf{0.8944} & 0.9217 & \textbf{0.9231} & $+$0.0015 & $+$0.0018 \\
GVM-CO     & 0.8887 & \textbf{0.8903} & 0.9188 & \textbf{0.9201} & $+$0.0016 & $+$0.0021 \\
LRE-CO     & 0.8948 & \textbf{0.8971} & 0.9240 & \textbf{0.9258} & $+$0.0023 & $+$0.0031 \\
LS-CO (G)  & 0.8916 & \textbf{0.8927} & 0.9213 & \textbf{0.9222} & $+$0.0011 & $+$0.0014 \\
LS-CO (C)  & 0.8918 & \textbf{0.8931} & 0.9199 & \textbf{0.9210} & $+$0.0013 & $+$0.0017 \\
\bottomrule
\end{tabular}
\end{table}

\begin{table}[H]
\centering
\caption{Average F1 improvement of no-PCA over PCA mode, stratified by dataset feature count.}
\label{tab:abl:pca:strat}
\begin{tabular}{lccccc}
\toprule
\textbf{Stratum} & \textbf{Datasets} & \textbf{PCA var.} & \multicolumn{3}{c}{$\boldsymbol{\Delta}$\textbf{F1 (no-PCA $-$ PCA)}} \\
 & & & GVM-CO & LRE-CO & LS-CO \\
\midrule
Low-dim ($d \leq 5$)       & haberman, new-thyroid1 & 92\% & $+$0.001 & $+$0.001 & $+$0.001 \\
Medium-dim ($6 \leq d \leq 9$) & ecoli, glass, yeast & 81\% & $+$0.002 & $+$0.003 & $+$0.001 \\
High-dim ($d \geq 13$)     & vehicle0, pima, heart, ionosphere & 61\% & $+$0.002 & $+$0.003 & $+$0.002 \\
\bottomrule
\end{tabular}
\end{table}

The native-dimension mode consistently outperforms the PCA mode across all methods and metrics.
The improvement is small in aggregate ($+0.001$--$+0.002$ F1) but meaningful at high IR ($+0.002$--$+0.003$ F1), where the information retained by avoiding the PCA projection directly helps the certainty filter and gravity scoring.
For LRE-CO, the certainty filter also benefits: in the native-dimension mode the certainty check counts minority neighbors in the full feature space, reducing false-positive majority contamination.

\textbf{Recommendation.}
The native-dimension mode (\texttt{use\_pca=False}) is recommended as the default for datasets with $d \leq 20$.
For very high-dimensional datasets ($d > 50$) or when runtime is critical, the PCA mode (\texttt{use\_pca=True}) remains appropriate.


% ============================================================================
\section{Execution Time Analysis}
\label{sec:abl:time}

\begin{table}[H]
\centering
\caption{Average execution time (seconds) per fold, measured on an Intel i7-12700H processor.}
\label{tab:abl:time}
\begin{tabular}{lccc}
\toprule
\textbf{Method} & \textbf{Small ($n < 300$)} & \textbf{Medium ($300 \leq n < 800$)} & \textbf{Large ($n \geq 800$)} \\
\midrule
None              & 0.001 & 0.002 & 0.005 \\
ROS               & 0.002 & 0.004 & 0.008 \\
SMOTE             & 0.010 & 0.025 & 0.060 \\
B-SMOTE           & 0.015 & 0.035 & 0.080 \\
ADASYN            & 0.012 & 0.030 & 0.070 \\
KM-SMOTE          & 0.020 & 0.050 & 0.120 \\
GVM-CO            & 0.025 & 0.065 & 0.180 \\
LRE-CO            & 0.080 & 0.220 & 0.650 \\
LS-CO ($L=3$)     & 0.015 & 0.040 & 0.100 \\
LS-CO ($L=10$)    & 0.040 & 0.110 & 0.290 \\
\bottomrule
\end{tabular}
\end{table}

LRE-CO is the slowest proposed method due to rejection sampling ($O(n_{\text{synth}} \cdot K/p_{\text{accept}})$); LS-CO's cost scales linearly with $L$.
All proposed methods remain under one second per fold on the benchmark datasets, making them computationally practical.
Total complexity:
\begin{itemize}[leftmargin=*]
    \item GVM-CO: $O(n_{\min} \log n_{\min} + n_{\text{synth}})$
    \item LRE-CO: $O(n_{\min} \log n_{\min} + n_{\text{synth}} / p_{\text{accept}} \cdot K)$
    \item LS-CO: $O(n_{\min} \log n_{\min} + n_{\text{synth}} \cdot L)$
\end{itemize}


% ============================================================================
\section{Chapter Summary}
\label{sec:abl:summary}

The ablation studies yield the following recommendations and findings:

\begin{itemize}[leftmargin=*]
    \item \textbf{Clustering method}: K-means outperforms HAC by 0.001--0.007 in F1 and is the default due to its alignment with circular geometry.
    \item \textbf{Denoising}: Tomek links provide a marginal benefit ($+0.001$--$+0.002$ F1); ENN degrades GVM-CO ($-0.019$) and LS-CO ($-0.015$) and should not be applied to these methods.
    \item \textbf{Hyperparameters}: The identified optima are $\alpha=0.4$ for GVM-CO (vs.\ default 0.6) and $L \geq 10$ for LS-CO (vs.\ default 3 chosen for runtime fairness).
    \item \textbf{Seed selection}: All four components (BC, JSD, AGTP, $Z$) contribute positively. AGTP and BC provide the largest individual contributions.
    \item \textbf{Progressive ablation}: The circular geometry produces the largest single change from SMOTE ($-0.003$ aggregate F1). Improvements are concentrated on high-IR datasets.
    \item \textbf{PCA for seed scoring}: No-PCA scoring ($k_{pc}=d$) is superior to 2-D PCA scoring ($k_{pc}=2$), improving AGTP by $+0.211$ and $Z$ by $-0.226$ on average.
    \item \textbf{PCA vs.\ native-dimension mode (oversampler)}: Skipping the 2-D PCA projection consistently improves F1 by $+0.001$--$+0.003$; recommended for $d \leq 20$.
    \item \textbf{Execution time}: All methods run under 1~s/fold. LRE-CO is slowest; LS-CO scales linearly with $L$.
\end{itemize}

\medskip\noindent\textbf{Bridge to Part~V.}
Part~IV has demonstrated that circular oversampling methods, particularly LRE-CO, improve minority-class classification at high imbalance ratios through circle-based generation, gravity-scoring, and BC-guided seed selection.
These improvements raise a deeper question: is the class imbalance that these methods are designed to correct the \emph{true} cause of poor minority-class performance?
Part~V steps back to ask this question directly through a controlled causal study of the imbalance--performance relationship, using the same GVM-CO algorithm developed in Part~IV as the recovery oversampler and random minority removal to isolate the effect of imbalance.
